\section{Introduction}
\label{sec:introduction}

A Banach-space norm defines both a linear topology and a metric.
A bounded linear isomorphism preserves the metric up to multiplicative
constants, but a map with that metric property need not preserve
addition or scalar multiplication. The Lipschitz-isomorphism problem
asks whether a \emph{bijection} with two-sided Lipschitz bounds nevertheless
forces two Banach spaces to be linearly isomorphic. We consider the
separable real case and obtain a negative answer.

Throughout, \(\ell_2\) denotes the real separable Hilbert space of
square-summable scalar sequences. For a Banach space \(H\), set
\[
 c_0(H)=\{(v_i)_{i\ge1}:v_i\in H,\ \|v_i\|\longrightarrow0\},
 \qquad \|(v_i)\|=\sup_i\|v_i\|.
\]
Ordinary \(c_0\) means \(c_0(\R)\). A linear embedding means a bounded
linear map that is bounded below; its range is consequently closed.

\begin{theorem}
\label{thm:main}
There are separable real Banach spaces \(X,Y\) and a bijection
\(\Psi:X\to Y\) such that, for every \(s,t\in X\),
\begin{equation}
\label{eq:main-bounds}
 \frac4{21}\|s-t\|_X
 \le \|\Psi(s)-\Psi(t)\|_Y
 \le \frac{76}{25}\|s-t\|_X.
\end{equation}
The space \(X\) contains a linearly isometric copy of \(c_0(\ell_2)\), while
\(Y\) contains no closed subspace linearly isomorphic to
\(c_0(\ell_2)\). In particular, \(X\) and \(Y\) are not linearly
isomorphic.
\end{theorem}

The obstruction in this theorem is deliberately block-valued:
each coordinate of \(c_0(\ell_2)\) is an entire Hilbert space.
It does not say that \(Y\) contains no ordinary \(c_0\).
The Hilbert block in each coordinate is essential to the obstruction
proved below.

\subsection*{Historical setting}

Ribe proved that uniformly homeomorphic normed spaces have the same
finite-dimensional linear structure up to a uniform distortion bound
\cite[Theorem~1]{Ribe76}. That result shows how much linear information
can survive even a weaker form of metric equivalence. At the same time,
Aharoni and Lindenstrauss constructed Lipschitz equivalent,
nonisomorphic nonseparable Banach spaces and explicitly asked for
separable examples \cite[pp.~281--282]{AL78}. Ribe later constructed
separable uniformly homeomorphic, nonisomorphic spaces \cite{Ribe84}.
Aharoni and Lindenstrauss then gave uniform pairs that are not
Lipschitz equivalent \cite{AL85}, while Johnson, Lindenstrauss and
Schechtman developed further separable uniform examples and left the
separable Lipschitz question open
\cite[Theorem~5.8 and Section~8(3)]{JLS96}. These developments make the
global two-sided estimate in \eqref{eq:main-bounds} a substantive
requirement, rather than a reformulation of a uniform counterexample.

Several positive results explain why the separable question resists a
direct transfer of those examples. Godefroy, Kalton and Lancien proved
that a Banach space Lipschitz isomorphic to \(c_0\) is linearly
isomorphic to \(c_0\), and that being linearly isomorphic to a subspace
of \(c_0\) is invariant under Lipschitz isomorphism
\cite[Theorems~2.1--2.2]{GKL00}. Godefroy and Kalton's lifting theorem
places every separable Banach space linearly and isometrically inside
its Lipschitz-free space; their quotient corollary converts a
Lipschitz right inverse of a linear quotient onto a separable space
into a bounded linear right inverse
\cite[Theorem~3.1 and Corollary~3.2]{GK03}. The construction below uses
free spaces; \Cref{rem:no-section} explains the distinction between
the lift used here and a right inverse of the linearized correction.

In the quasi-Banach setting, Albiac and Kalton obtained separable
counterexamples for every \(0<p<1\) \cite[Theorem~5.4]{AK09}.
The strict inequality \(p<1\) is part of that theorem, so it does not
settle the Banach-norm case. Kalton recorded the separable
Lipschitz-isomorphism question as Problem~3 in his 2008 survey
\cite[Section~2.2]{Kalton08}; Albiac, Ansorena, B\'{i}ma and C\'{u}th
still described the Banach case as open in the first paragraph of
their paper published online in December 2025 \cite{AABC26}.
Theorem~\ref{thm:main} resolves that question negatively for separable
real Banach spaces. It asserts neither reflexivity nor a complex-scalar
analogue.

For comparison, every surjective isometry between the unit spheres of
nonzero real Banach spaces extends uniquely to a surjective real-linear
isometry \cite[Theorem~1.1]{OpenAITingley}.
That result concerns exact distances on unit spheres; here we allow
fixed multiplicative distortion but require a bijection of the whole
spaces.

\subsection*{Proof overview}

The proof separates the metric construction from the linear
obstruction. Its intermediate map acts between Hilbert spaces:
\[
 M=U\oplus_2 V,\qquad V=\ell_2,\qquad
 h:M\longrightarrow U.
\]
Here \(U=(\bigoplus_{n\ge1}U_n)_2\) is a Hilbert sum of
finite-dimensional coordinate blocks. The map \(h\) is an onto
bi-Lipschitz change of variables fixing zero, close to an isometry. Its relevant
nonlinear part is
\[
 D(u,v)=h(u,v)-u.
\]
The construction arranges that, near any finite family of distinct
points of \(M\), the sufficiently high output coordinates of \(D\)
lie in mutually orthogonal coordinate subspaces. The cutoff index and
the neighborhoods may depend on that finite family. We call this
property \emph{local orthogonality}.

To turn this geometry into a linear obstruction, use the
Lipschitz-free space \(E=\F(M)\), the closed linear span of evaluation
functionals on real zero-preserving Lipschitz functions.
The evaluation map \(\delta:M\to E\) preserves distances. Every
zero-preserving Lipschitz map has a bounded linearization, so \(D\)
determines \(q:E\to U\) by
\[
 q\delta(u,v)=h(u,v)-u.
\]
Section~\ref{sec:localization} proves that local orthogonality makes
\(q\) \emph{completely continuous}: a weakly null sequence in \(E\)
has norm-null image. This is the one place where the local coordinate
geometry enters the Banach-space obstruction.

The linearization is also what makes the metric change of variables
possible. Give \(E\oplus V\) its sum norm and put
\[
 Q_1(e,v)=(qe,v),\qquad
 B(e,v)=e+\delta(qe,v).
\]
The identity \(qB=hQ_1\) is triangular: applying \(q\) to the new
variable records \(h\) of the old output. Given \(z\in E\), one first
recovers \(x=h^{-1}(qz)\) and then subtracts \(\delta(x)\); this yields
the explicit inverse of \(B\). Section~\ref{sec:triangular} proves
both inverse identities and keeps track of the sum norm.

For a bounded linear \(Q:A\to H\), the graph space used in the final
step consists of sequences \(s=(s_i)\) for which
\(\sum_i2^{-i}\|s_i\|<\infty\) and \((Qs_i)\in c_0(H)\), with norm
\[
 \|s\|_{Z_Q}
 =\max\left\{\sum_i2^{-i}\|s_i\|,\ \sup_i\|Qs_i\|\right\}.
\]
Set \(X=Z_{Q_1}\) and \(Y=Z_q\). Applying \(B\) in each coordinate
gives the desired bijection: \(qB=hQ_1\) controls the second part of
the graph norm, and the Lipschitz bounds for \(B\) control the weighted
sum. The \(V\) coordinates embed \(c_0(V)\) isometrically in \(X\).
By contrast, complete continuity of \(q\) makes each fixed output
coordinate tend to zero along an orthonormal sequence in any Hilbert
block of a proposed \(c_0(\ell_2)\) copy in \(Y\).
A disjoint-support argument then finds a unit vector in each block
with arbitrarily small full \(c_0(U)\) output. A surviving copy would embed
ordinary \(c_0\) into \(\ell_1(E)\). Section~\ref{sec:free-graph}
excludes this because \(E=\F(M)\) is weakly sequentially complete.

Two analytic steps supply the inputs just described. For localized
linearization, a hypothetical non-null image of a weakly null
sequence produces uniformly Lipschitz scalar tests \(d_i\).
Orthogonality allows only finitely many points where their limiting
local Lipschitz constants remain large. Cutoffs remove neighborhoods
of those points; a shrinking-ball lemma shows that the removed tests
pair negligibly with the weakly null sequence. The compact reduction
theorem of Aliaga, No\^us, Petitjean and Proch\'azka
\cite[Theorem~2.3]{ANPP21} then reduces the remainder to a compact set,
where a finite cover controls its Lipschitz constant.

For the Hilbert map, a unit vector in each block is designated as a
slot. With the slots frozen, perpendicular components stay in their
blocks, old slot coordinates move to even blocks, and the additional
\(\ell_2\) coordinates fill odd blocks. This is an onto linear
isometry. The slots are then allowed to depend on the input so that
near different points they eventually use disjoint coordinates.
Finite approximants freeze all but finitely many slots and have
uniform global inverse bounds. Section~\ref{sec:hilbert} constructs the slots with a common radial
derivative budget, proves global invertibility of the finite approximants,
and passes to an onto limit. Its decisive estimate bounds every tail of the
inverse images by the corresponding tail of a fixed output, with an
exact cancellation of slot terms. Section~\ref{sec:assembly} combines
the results and computes the constants in \eqref{eq:main-bounds}.
